\section{The leading interaction of the two boundaries}
\label{v13:sec:boundary-interaction}

We evaluate the full interaction, including the complementary-space
correction, which has the same order as the direct coupling. The
reduction is analogous to Feshbach--Schur
\cite[Theorem~1.2 and equations~(1.15)--(1.16)]{DussonSigalStamm2021};
the physical sine-Loewner modes determine its explicit coefficient.

Fix a real $\alpha\in(-g,0)\cup(0,g)$, and write
\begin{equation}\label{v13:eq:interaction-parameters}
 a=\sqrt{1-2\alpha^2},\qquad
 \beta=\pi^{-1}\arctan a,\qquad \nu=\tfrac12+\beta,\qquad
 \gamma=2\beta,\qquad
 c_\alpha=-\frac{\pi a(1+a^2)}{\alpha}.
\end{equation}
For a normalized left boundary eigenfunction $\psi_L$ on $(l,\infty)$,
define its physical tail amplitude by
\begin{equation}\label{v13:eq:physical-amplitude}
 \eta_L=\lim_{T\to\infty}T^\nu
                 \int_T^{T+1}\psi_L(l+x)\,dx
 =\frac{\Gamma(\nu)}{(2\pi)^\nu}
       \left(e^{-i\pi\nu/2}v_1+e^{i\pi\nu/2}v_4\right),
\end{equation}
where $v$ is the endpoint vector in
Theorem~\ref{v12:thm:boundary-tail}, scaled so that its physical
eigenfunction has norm one on the half-line.  The limit follows from
\eqref{v12:eq:physical-tail}: integration over one unit interval
removes the two nonzero leading carrier frequencies.  For a normalized
right boundary eigenfunction $\psi_R$ on $(-\infty,h)$, define
$\eta_R$ by the same limit with $\psi_R(h-x)$ in the integrand.
The amplitudes may vanish; their phase follows the eigenfunction's phase.
The same amplitude is determined in modulus by the normalized physical
tail mass:
\[
 \lim_{T\to\infty}T^\gamma\int_T^\infty
       |\psi_L(l+x)|^2\,dx=\frac{|\eta_L|^2}{\alpha^2\gamma}.
\]
Indeed the tail direction in the proof has squared norm $\alpha^{-2}$;
the nonmatching carriers integrate to a lower order. The right analogue
holds with $\psi_R(h-x)$.

\begin{theorem}[Closed leading boundary interaction]
\label{v13:thm:leading-interaction}
Fix $l,h$, and suppose that $\alpha$ is an isolated eigenvalue of one
or both half-line operators in Theorem~\ref{v11:thm:gap-limits}.
Their individual eigenvalues are simple by
Theorem~\ref{boundary:thm:one-cell-flow}.  Set
\[
 T_m=P_{(l,m+h)}\mathcal G P_{(l,m+h)},\qquad L=m+h-l,
 \qquad m\longrightarrow\infty\quad(m\in\mathbb Z).
\]
If $\alpha$ belongs to both boundary spectra, let $\mu_1\le\mu_2$
be the eigenvalues of the Hermitian matrix
\begin{equation}\label{v13:eq:closed-interaction-matrix}
 \mathsf M_\alpha=c_\alpha
 \begin{pmatrix}0&\overline{\eta_L}\eta_R\\
                 \overline{\eta_R}\eta_L&0\end{pmatrix}.
\end{equation}
The two finite-interval eigenvalues converging to $\alpha$, in
increasing order within the cluster, satisfy
\begin{equation}\label{v13:eq:leading-cluster-law}
 \lambda_j(T_m)=\alpha+L^{-\gamma}\mu_j+o(L^{-\gamma}),
 \qquad j=1,2.
\end{equation}
If $\alpha$ belongs to only one boundary spectrum, its single
eigenvalue satisfies $\lambda(T_m)=\alpha+o(L^{-\gamma})$.
All parameters are fixed; no uniformity as $\alpha$ approaches zero
or a band edge is asserted.
\end{theorem}

The two-boundary case implies $l+h\in\mathbb Z$.
Indeed reflection identifies the right operator at $h$ with the left
operator at $-h$, and Theorem~\ref{boundary:thm:one-cell-flow}
assigns a fixed nonzero gap level to a unique boundary phase modulo one.
Conversely, this congruence makes the two half-line operators unitarily
equivalent. Centered intervals satisfy the congruence automatically.

\begin{corollary}[Centered parity splitting]
\label{v13:cor:leading-parity-splitting}
Let $\alpha$ be an isolated eigenvalue of $L_{-r}\mathcal G L_{-r}$,
with normalized eigenfunction of amplitude $\eta$ in
\eqref{v13:eq:physical-amplitude}.  For the centered operator
$\mathcal K_{n+r}$, put $L=2(n+r)$.  Its unique even and odd
eigenvalues converging to $\alpha$ satisfy
\begin{align}
 \lambda_{\rm even}
   &=\alpha+c_\alpha|\eta|^2L^{-\gamma}+o(L^{-\gamma}),\nonumber\\
 \lambda_{\rm odd}
   &=\alpha-c_\alpha|\eta|^2L^{-\gamma}+o(L^{-\gamma}).
 \label{v13:eq:closed-parity-law}
\end{align}
Their midpoint is $\alpha+o(L^{-\gamma})$.  If $\eta\ne0$, their
splitting has the nonzero leading magnitude
$2|c_\alpha||\eta|^2L^{-\gamma}$, and eventually
$|\lambda_{\rm even}|<|\lambda_{\rm odd}|$.
If $\eta=0$, both shifts are $o(L^{-\gamma})$; strict ordering is
not inferred.  In fact both shifts are $O(L^{-1+\delta})$ for every
fixed $\delta>0$ in this case.
\end{corollary}

Section~\ref{v13:app:boundary-interaction} proves the full correction
and evaluates its scalar integral. Theorem~\ref{boundary:thm:certified-residue:rr}
certifies $\eta\ne0$ for the phase-zero mode.

